% Supplementary Material -- Paper 6 (T6), IEEE JBHI special issue submission.
% Tables are generated by code/supplement_tables.py (and code/demand_table.py) from results/analysis-ablation-2026-10-07/
% and the CFD returns; do not edit the numbers by hand.
\documentclass[journal,onecolumn,11pt]{IEEEtran}
\usepackage{amsmath,amssymb,booktabs,siunitx}
\usepackage[hidelinks]{hyperref}
\usepackage{placeins}
\renewcommand{\thetable}{S\arabic{table}}
\renewcommand{\thesection}{S\arabic{section}}
\makeatletter\let\inputtab\@@input\makeatother

\begin{document}
\sloppy
\title{Supplementary Material for ``Topological Segmentation Error and Boundary-Condition Tuning in Computed Coronary
FFR: A Controlled In Silico Study''}
\author{Fadillah Yamin, Mohd Azan Mohammed Sapardi, Ming Kwang Tan, Xin Wang, and Mohd-Zulhilmi Paiz Ismadi}
\maketitle

\section{Departures From the Analysis Plan}
The cohort, design, analysis plan and ablation code were fixed, and the cohort list hashed, before the ablation was
run; the analysis script was written while the ablation ran and before its output was read. The following departures
from the plan occurred. None changes a conclusion stated in the main text.
\begin{enumerate}
\item \textit{Mixed-effects model.} The plan specified a mixed-effects logistic regression for flips. No frequentist
logistic mixed model was available in the analysis environment, so the same model (error type, protocol and their
interaction, severity, length and baseline band as fixed effects; random intercepts for patient and lesion slot) was
fitted as a Bayesian model by variational inference (Table~\ref{tab:glmm}). A patient-clustered generalized
estimating equation, fitted as an additional check that was not in the plan, gave estimates in the discrete bed with
the same sign as the Bayesian model for every error-type and protocol term, but a degenerate coefficient for the
0.90--0.95 band, in which no model flipped. In the leaky bed, where no model in the two upper bands flipped, it did not
converge. The main-text inferences rest on the paired McNemar, sign and Wilcoxon tests.
\item \textit{Paired $\Delta$FFR model.} The plan specified a paired mixed model for $\Delta$FFR with the same terms as
the flip model. It was not fitted; paired changes in $|\Delta\mathrm{FFR}|$ were tested with the Wilcoxon signed-rank
test.
\item \textit{Bed structure.} The plan specified a pooled model with bed as a fixed effect and an error-type $\times$
bed interaction. The beds were analyzed separately, and a finding is claimed only where its direction agrees in both.
\item \textit{Multiplicity.} The plan specified Holm adjustment across all hypotheses. The Holm adjustment was applied
to the three protocol contrasts within each error type and bed; the remaining tests are reported unadjusted.
\item \textit{Discrete-bed eligibility.} The ablation was run on all 150 instances in the leaky bed and on the 148
whose lesion slot could be built in the discrete bed. The rule that defines the discrete arm (healthy-equivalent
network FFR $\geq 0.90$ and at least two outlets; 97 instances, all among the 148) was applied at the analysis stage.
\item \textit{Per-territory tuning.} The plan specified a sensitivity analysis in which Protocol C fits one scaling
per territory. It was not implemented and was not run.
\item \textit{Noise floor.} The repeat-measurement floor was computed per instance as the probability that a normal
deviation with standard deviation 0.018 about the clean FFR crosses 0.80, rather than read from published
classification curves. The simulated floor used 20 draws per instance; the number of draws was not specified in the plan.
\item \textit{Demand sensitivity.} The plan specified a sensitivity to the hyperemic demand scaled by 0.7 and 1.3. It
was run on the ablation after the primary analysis (Section~\ref{sec:demand}).
\item \textit{Secondary analyses.} The plan listed stratified analyses by branch position relative to the lesion,
bifurcation lesion, image quality, host vessel and gray zone (FFR 0.75--0.85). They were not run. Per-vessel counts
were tabulated but are not reported as a stratum, because Protocol C is undefined for most right coronary instances
with a topological error (main text, Limitations).
\item \textit{Instances solved before the run.} Eight instances had been solved during code verification and in the
three-dimensional case before the ablation was run. All eight remain in the cohort; the planned sensitivity analysis
excluding them was not run.
\item \textit{Uncertainty floors.} Beside the reduced-order discretization floor of 0.005, the plan listed a
healthy-reference re-fit floor of 0.010 and a tolerance of $+3\%$ on the calibrated healthy-equivalent inflow. No FFR
difference below 0.01 is interpreted as meaningful unless it is a paired within-instance contrast.
\item \textit{Scope.} The plan also specified a 0D--3D agreement analysis on 30 instances and a deployment-time error
detector. Their results are not part of this report.
\item \textit{Three-dimensional case.} Three mesh acceptance rules were met only with declared deviations, and the
throat-resolution test did not meet its acceptance criterion; both are reported in Section~\ref{sec:cfd}, with the
resolution sensitivity stated beside the error effect.
\end{enumerate}

\FloatBarrier
\section{Exclusions}
Of 6\,944 swept instances, 150 were selected and 97 form the discrete arm; 2\,856 corrupted models were attempted and
2\,599 solved, all of which converged. The missed branch was applicable to 77 of 97 (discrete) and 118 of 150 (leaky)
instances; the others had no side branch distal to the lesion that could be deleted. The vessel break was not
applicable to one instance in each bed, in which the 25~mm truncation point fell at the start of a centerline segment.
Protocol C is defined only when at least two perfusion territories retain a vessel. A fit that ends at a search bound
is recorded as a failure; none did in the ablation or in the demand runs (Section~\ref{sec:demand}), and 10 of
5\,860 noise-floor draws did (Section~\ref{sec:floor}). Under Protocol A a vessel break can remove every node that
carries bed outflow (36 instances in the discrete bed, 2 in the leaky bed). Table~\ref{tab:excl} lists every
exclusion; all other cells were solved for every applicable instance.
\begin{table}[!ht]
\caption{Cells With Excluded Models}\label{tab:excl}
\centering\small
\begin{tabular}{@{}lllrl@{}}
\toprule
Bed & Error & Protocol & Solved & Excluded (reason: count)\\
\midrule
\inputtab supplement_tables/tab_exclusions_nz.tex 
\bottomrule
\end{tabular}
\end{table}

\FloatBarrier
\section{Mixed-Effects Model of Flips}
\begin{table}[!ht]
\caption{Bayesian Mixed-Effects Logistic Model of Flips: Posterior Mean Log-Odds (95\% Interval)}\label{tab:glmm}
\centering\small
\begin{tabular}{@{}lcc@{}}
\toprule
Term & Discrete bed & Leaky bed\\
\midrule
\inputtab supplement_tables/tab_glmm.tex 
\bottomrule
\end{tabular}
\\[3pt]
\parbox{0.85\textwidth}{\footnotesize Reference levels: T1 (missed branch) and Protocol A. Baseline-band terms are
omitted from the table. Random intercepts for patient and lesion slot. Fitted by variational inference; intervals are
posterior mean $\pm 1.96$ posterior standard deviations from the variational approximation.}
\end{table}

At the reference levels the model agrees with the paired tests: under fixed boundary conditions a vessel break carries
higher odds of a flip than a missed branch and the caliber errors lower odds, and for the missed branch Protocols B and
C carry lower odds than Protocol A, in both beds. In the leaky bed the T2 $\times$ B and T2 $\times$ C interactions
offset the T2 main effect, so under re-derived or tuned boundary conditions the two topological errors have similar
odds, as in Table I of the main text.

\FloatBarrier
\section{Sensitivity Analyses}
\subsection{Validation Threshold}
Table~\ref{tab:thr} repeats the proportion of models that pass the perfusion check while materially wrong at the
primary threshold of 10\% and at the 13\% and 16\% sensitivities, for the topological errors and for all error
types.
\begin{table}[!ht]
\caption{Proportion of Models That Pass the Perfusion Check and Are Materially Wrong, \% (Wilson 95\% Interval), by
Validation Threshold}\label{tab:thr}
\centering\small
\begin{tabular}{@{}lllrccc@{}}
\toprule
Bed & Errors & Protocol & $n$ & 10\% & 13\% & 16\%\\
\midrule
\inputtab supplement_tables/tab_thresholds.tex 
\bottomrule
\end{tabular}
\\[3pt]
\parbox{0.85\textwidth}{\footnotesize Materially wrong: $|\Delta\mathrm{FFR}| > 0.05$. The 13\% and 16\% thresholds were analyzed as sensitivities. $n$ counts models with a
defined perfusion residual: T2 under Protocols A and B in the leaky bed, $n = 100$; T2 under Protocol B in the
discrete bed, $n = 60$.}
\end{table}

\FloatBarrier
\subsection{Hyperemic Demand}\label{sec:demand}
The demand constant $k$ was scaled by 0.7 and 1.3 and the full ablation was repeated with the same frozen cohort
(Table~\ref{tab:demand}). The ordering of topological against caliber errors under fixed boundary conditions, and of
Protocols B and C against Protocol A for the topological errors, was unchanged at both levels in both beds. The taper,
which flipped fewer decisions under Protocol C than under Protocol A at the primary demand, flipped one and two more
at 1.3 times the demand (7 against 6 of 97; 11 against 9 of 150). The magnitudes changed with demand, which moves the
clean FFR of every instance (leaky-bed median 0.869, 0.801 and 0.741 at 0.7, 1.0 and 1.3 times the demand) and
therefore the number of instances near the threshold, as well as the pressure loss across the lesion. The proportion
of tuned topological-error models that pass the perfusion check while materially wrong fell to 3\% in the leaky bed
at 0.7 times the demand, comparable with the simulated floor for correct anatomy (3.9\% at the primary demand); in
the discrete bed it remained between 12\% and 25\%.
\begin{table}[!h]
\caption{Sensitivity to Hyperemic Demand}\label{tab:demand}
\centering\small
\begin{tabular}{@{}llcccccc@{}}
\toprule
Bed & Demand scale & \multicolumn{2}{c}{Flip, \% (Protocol A)} & \multicolumn{2}{c}{Topological flip, \%} &
Passes and wrong, \% & Missed-branch $\Delta$FFR\\
\cmidrule(lr){3-4}\cmidrule(lr){5-6}
 & & Topological & Caliber & B & C & (C, topological) & median, A / C\\
\midrule
Discrete & 0.7 & 32 (25--40) & 6 (3--10) & 12 (8--18) & 18 (12--27) & 12 (7--19) & 0.070 / 0.057 \\
 & 1.0 & 33 (26--41) & 10 (7--15) & 14 (10--20) & 19 (13--28) & 19 (13--28) & 0.095 / 0.077 \\
 & 1.3 & 40 (32--48) & 5 (2--9) & 16 (11--22) & 19 (13--28) & 25 (18--34) & 0.106 / 0.091 \\
Leaky & 0.7 & 20 (16--26) & 6 (4--9) & 7 (4--10) & 9 (5--14) & 3 (1--7) & 0.032 / 0.010 \\
 & 1.0 & 32 (26--38) & 6 (4--9) & 5 (3--8) & 8 (5--13) & 8 (5--13) & 0.047 / 0.015 \\
 & 1.3 & 37 (32--43) & 4 (2--7) & 5 (3--9) & 8 (4--13) & 12 (8--18) & 0.057 / 0.021 \\
\bottomrule
\end{tabular}
\\[3pt]
\parbox{0.9\textwidth}{\footnotesize Demand scale 1.0 is the primary analysis ($k = 562$~s$^{-1}$). Percentages with Wilson 95\%
intervals in parentheses. Missed-branch $\Delta$FFR is over the instances in which Protocol C was defined.}
\end{table}


\FloatBarrier
\section{Noise Floors}\label{sec:floor}
The repeat-measurement floor uses the published standard deviation of repeat invasive FFR, 0.018. The simulated floor
uses the clean anatomy of every instance, with 20 draws per instance. Each draw perturbs cardiac output (relative
within-subject standard deviation 0.05), mean aortic pressure (0.056), blood viscosity through hematocrit (0.02), the
share of flow to each territory (0.10) and the measured territory flows (within-subject coefficient of variation
0.083), and the model is then tuned to the perturbed flows as in Protocol C. Heart rate is not represented in the
steady model. The territory-share and measurement terms lie at the low end of published ranges, so the simulated floor
may understate physiological variation. Table~\ref{tab:floor} gives the result. The repeat-measurement floor is
5.0--6.5\% across cells and the simulated floor 6.2\% and 7.2\%.
\begin{table}[!ht]
\caption{Simulated Floor: Correct Anatomy Tuned to Noisy Perfusion}\label{tab:floor}
\centering\small
\begin{tabular}{@{}lrrcccrc@{}}
\toprule
Bed & Instances & Draws solved & Flip, \% & Median $|\Delta$FFR$|$ & 95th pct. $|\Delta$FFR$|$ & Pass, \% &
Pass and wrong, \%\\
\midrule
\inputtab supplement_tables/tab_floor.tex 
\bottomrule
\end{tabular}
\end{table}

\FloatBarrier
\section{Three-Dimensional Case Study: Mesh Checks}\label{sec:cfd}
The mask-to-solution path was scripted with no manual repair, and the clean, baseline and missed-branch solves
converged. Three acceptance rules were met with deviations. (i) The throat-size gate compared the ratio of the
as-built throat radius to the undeformed meshed radius with the intended radial scale (within 1\%), rather than the
absolute throat radius with the reduced-order target, because the two pipelines define the radius differently. (ii) A surface self-intersection flag arises from the segmentation mask, not from the mesher. (iii) The
strict mesh-quality check flags poorly shaped cells (face-tetrahedra decompositions) 1.43~mm from the throat center,
inside the planned 2~mm exclusion distance; the lesion-free mesh carries a similar count, and the standard quality
check passes. To test whether these cells affect the result, the baseline was re-meshed with a finer throat zone
(Table~\ref{tab:d7}). The acceptance criterion, a change smaller than the discretization uncertainty of
$5.5\times10^{-4}$, was not met: the finer zone raised the measurement-point FFR by 0.0019 and 0.0021, while a
regenerated mesh at the production resolution reproduced the original value. The second variant missed the strict
residual limit at the iteration budget (final velocity residual $1.08\times10^{-5}$ against $10^{-5}$) with its inlet
pressure steady to $4\times10^{-4}$~Pa over the last 500 iterations; its value is reported as returned. This is a resolution sensitivity of
about 0.2\% of the FFR, against an error effect of 0.074 in the same case.

All 3D-to-0D comparisons use a reduced-order twin rebuilt on the meshed radius, which is wider than the
reduced-order radius (main text, Section III-D).
\begin{table}[!ht]
\caption{Throat-Zone Resolution Test, Baseline Lumen. The 8.31~M-cell mesh did not meet the strict residual limit at
3\,000 iterations (see text).}\label{tab:d7}
\centering\small
\begin{tabular}{@{}lrrrrr@{}}
\toprule
Mesh & Cells (M) & FFR & $\Delta$FFR & Outlet flow change, \% & Flagged cell to throat, mm\\
\midrule
\inputtab supplement_tables/tab_d7.tex 
\bottomrule
\end{tabular}
\end{table}

\end{document}
